% GENERATED FILE. Edit canonical lessons, then run npm run generate:books.
% canonical-source-hash: fnv1a64:392147fafcffcecb
% canonical-lessons: JA-C228-manabu, JA-C228-yakusu, JA-C228-namakeru, JA-C228-yameru, JA-C228-sotsugyou

\chapter{Learning, Slacking and Graduating}
\label{ch:ja-learning-slacking-and-graduating}
\hlchaptermodality{228}

\begin{chapteropening}
\textbf{By the end of this chapter:} \emph{I can talk about learning, quitting and finishing school.}

Complete the last lesson of chapter 228: I can talk about learning, quitting and finishing school.
\end{chapteropening}

\section[\texorpdfstring{\ja{まなぶ} (manabu)}{まなぶ (manabu)}]{\ja{まなぶ} (manabu) — to learn; to study in depth}

\label{lesson:JA-C228-manabu}



\begin{quote}
Before the new one: say the Japanese for to make sure, then the Japanese for to compare.
\end{quote}

\subsection*{You'll want to know: \ja{まなぶ}}
\textbf{\ja{まなぶ}} — \emph{manabu} — \textquotedblleft{}to learn; to study in depth\textquotedblright{}.

\begin{quote}
\textbf{\ja{あねは} \ja{がっこうで} \ja{にほんごを} \ja{まなびます。}} — \emph{ane wa gakkō de nihongo o manabimasu} — \textquotedblleft{}My elder sister studies Japanese at school.\textquotedblright{}
\end{quote}

Polite \textbf{\ja{まなびます}}, \emph{manabimasu}; a godan verb.

\begin{cousinweb}
\textbf{\ja{まなぶ}} (\emph{manabu}) is a native Japanese word (wago), inherited rather than borrowed from Chinese; its kanji was attached later as a kun'yomi (JMdict, KANJIDIC2).
\end{cousinweb}

\begin{grammarlens}[title={What you've built}]
One more word for everyday life at A2.
\end{grammarlens}

\subsection*{Your turn}
\begin{itemize}
\raggedright
  \item \emph{Say it:} \emph{manabu}
  \item \emph{Say it:} \emph{manabu}, once more
  \item \emph{Say it:} \emph{ane wa gakkō de nihongo o manabimasu} — My elder sister studies Japanese at school.
  \item \emph{Recall:} say the Japanese for to make sure, then the Japanese for to compare, then say \emph{manabu} again
\end{itemize}

\begin{tcolorbox}[breakable,colback=teal!4,colframe=teal!35!black,title={Before you move on}]
What does \ja{まなぶ} mean? (\textquotedblleft{}To learn; to study in depth\textquotedblright{}.) Say it once more.
\end{tcolorbox}
\section[\texorpdfstring{\ja{やくす} (yakusu)}{やくす (yakusu)}]{\ja{やくす} (yakusu) — to translate}

\label{lesson:JA-C228-yakusu}



\begin{quote}
Before the new one: say the Japanese for to compare, then the Japanese for to learn.
\end{quote}

\subsection*{You'll want to know: \ja{やくす}}
\textbf{\ja{やくす}} — \emph{yakusu} — \textquotedblleft{}to translate\textquotedblright{}.

\begin{quote}
\textbf{\ja{ちちは} \ja{にほんごの} \ja{てがみを} \ja{やくします。}} — \emph{chichi wa nihongo no tegami o yakushimasu} — \textquotedblleft{}My father translates letters written in Japanese.\textquotedblright{}
\end{quote}

Polite \textbf{\ja{やくします}}, \emph{yakushimasu}; a godan verb.

\begin{cousinweb}
\textbf{\ja{やくす}} (\emph{yakusu}) is a mixed word, built on Sino-Japanese \emph{yaku}, \textquotedblleft{}translate\textquotedblright{} (Middle Chinese \emph{yek}), borrowed from Chinese, with a native verb ending \emph{-su} (JMdict, KANJIDIC2; Middle Chinese in Baxter's transcription).
\end{cousinweb}

\begin{grammarlens}[title={What you've built}]
One more word for everyday life at A2.
\end{grammarlens}

\subsection*{Your turn}
\begin{itemize}
\raggedright
  \item \emph{Say it:} \emph{yakusu}
  \item \emph{Say it:} \emph{yakusu}, once more
  \item \emph{Say it:} \emph{chichi wa nihongo no tegami o yakushimasu} — My father translates letters written in Japanese.
  \item \emph{Recall:} say the Japanese for to compare, then the Japanese for to learn, then say \emph{yakusu} again
\end{itemize}

\begin{tcolorbox}[breakable,colback=teal!4,colframe=teal!35!black,title={Before you move on}]
What does \ja{やくす} mean? (\textquotedblleft{}To translate\textquotedblright{}.) Say it once more.
\end{tcolorbox}
\section[\texorpdfstring{\ja{なまける} (namakeru)}{なまける (namakeru)}]{\ja{なまける} (namakeru) — to slack off; to be idle}

\label{lesson:JA-C228-namakeru}



\begin{quote}
Before the new one: say the Japanese for to learn, then the Japanese for to translate.
\end{quote}

\subsection*{You'll want to know: \ja{なまける}}
\textbf{\ja{なまける}} — \emph{namakeru} — \textquotedblleft{}to slack off; to be idle\textquotedblright{}.

\begin{quote}
\textbf{\ja{おとうとは} \ja{よく} \ja{しゅくだいを} \ja{なまけます。}} — \emph{otōto wa yoku shukudai o namakemasu} — \textquotedblleft{}My younger brother often neglects his homework.\textquotedblright{}
\end{quote}

Polite \textbf{\ja{なまけます}}, \emph{namakemasu}; an ichidan verb.

\begin{cousinweb}
\textbf{\ja{なまける}} (\emph{namakeru}) is a native Japanese word (wago), inherited rather than borrowed from Chinese; its kanji was attached later as a kun'yomi (JMdict, KANJIDIC2).
\end{cousinweb}

\begin{grammarlens}[title={What you've built}]
One more word for everyday life at A2.
\end{grammarlens}

\subsection*{Your turn}
\begin{itemize}
\raggedright
  \item \emph{Say it:} \emph{namakeru}
  \item \emph{Say it:} \emph{namakeru}, once more
  \item \emph{Say it:} \emph{otōto wa yoku shukudai o namakemasu} — My younger brother often neglects his homework.
  \item \emph{Recall:} say the Japanese for to learn, then the Japanese for to translate, then say \emph{namakeru} again
\end{itemize}

\begin{tcolorbox}[breakable,colback=teal!4,colframe=teal!35!black,title={Before you move on}]
What does \ja{なまける} mean? (\textquotedblleft{}To slack off; to be idle\textquotedblright{}.) Say it once more.
\end{tcolorbox}
\section[\texorpdfstring{\ja{やめる} (yameru)}{やめる (yameru)}]{\ja{やめる} (yameru) — to quit; to resign}

\label{lesson:JA-C228-yameru}



\begin{quote}
Before the new one: say the Japanese for to translate, then the Japanese for to slack off.
\end{quote}

\subsection*{You'll want to know: \ja{やめる}}
\textbf{\ja{やめる}} — \emph{yameru} — \textquotedblleft{}to quit; to resign\textquotedblright{}.

\begin{quote}
\textbf{\ja{あには} \ja{らいげつ} \ja{しごとを} \ja{やめます。}} — \emph{ani wa raigetsu shigoto o yamemasu} — \textquotedblleft{}My elder brother quits his job next month.\textquotedblright{}
\end{quote}

Polite \textbf{\ja{やめます}}, \emph{yamemasu}; an ichidan verb.

\begin{cousinweb}
\textbf{\ja{やめる}} (\emph{yameru}) is a native Japanese word (wago), inherited rather than borrowed from Chinese; its kanji was attached later as a kun'yomi (JMdict, KANJIDIC2).
\end{cousinweb}

\begin{grammarlens}[title={What you've built}]
One more word for everyday life at A2.
\end{grammarlens}

\subsection*{Your turn}
\begin{itemize}
\raggedright
  \item \emph{Say it:} \emph{yameru}
  \item \emph{Say it:} \emph{yameru}, once more
  \item \emph{Say it:} \emph{ani wa raigetsu shigoto o yamemasu} — My elder brother quits his job next month.
  \item \emph{Recall:} say the Japanese for to translate, then the Japanese for to slack off, then say \emph{yameru} again
\end{itemize}

\begin{tcolorbox}[breakable,colback=teal!4,colframe=teal!35!black,title={Before you move on}]
What does \ja{やめる} mean? (\textquotedblleft{}To quit; to resign\textquotedblright{}.) Say it once more.
\end{tcolorbox}
\section[\texorpdfstring{\ja{そつぎょう} (sotsugyō)}{そつぎょう (sotsugyō)}]{\ja{そつぎょう} (sotsugyō) — graduation}

\label{lesson:JA-C228-sotsugyou}



\begin{quote}
Before the new one: say the Japanese for to slack off, then the Japanese for to quit.
\end{quote}

\subsection*{You'll want to know: \ja{そつぎょう}}
\textbf{\ja{そつぎょう}} — \emph{sotsugyō} — \textquotedblleft{}graduation\textquotedblright{}.

\begin{quote}
\textbf{\ja{あねは} \ja{らいねん} \ja{そつぎょうです。}} — \emph{ane wa rainen sotsugyō desu} — \textquotedblleft{}My elder sister graduates next year.\textquotedblright{}
\end{quote}

\begin{cousinweb}
\textbf{\ja{そつぎょう}} (\emph{sotsugyō}) is Sino-Japanese (kango): two Chinese characters read with their on'yomi, readings borrowed from Chinese. They are \emph{sotsu}, \textquotedblleft{}graduate, finish\textquotedblright{} (Middle Chinese \emph{tswit}) and \emph{gyō}, \textquotedblleft{}work, vocation\textquotedblright{} (\emph{ngjaep}) (KANJIDIC2; Middle Chinese in Baxter's transcription).
\end{cousinweb}

\begin{grammarlens}[title={What you've built}]
One more word for everyday life at A2.
\end{grammarlens}

\subsection*{Your turn}
\begin{itemize}
\raggedright
  \item \emph{Say it:} \emph{sotsugyō}
  \item \emph{Say it:} \emph{sotsugyō}, once more
  \item \emph{Say it:} \emph{ane wa rainen sotsugyō desu} — My elder sister graduates next year.
  \item \emph{Recall:} say the Japanese for to slack off, then the Japanese for to quit, then say \emph{sotsugyō} again
\end{itemize}

\begin{tcolorbox}[breakable,colback=teal!4,colframe=teal!35!black,title={Before you move on}]
What does \ja{そつぎょう} mean? (\textquotedblleft{}Graduation\textquotedblright{}.) Say it once more.
\end{tcolorbox}
